\chapter{Chain Complexes}
\label{ch:viii16-chain-complexes}

Homology begins with a simple algebraic pattern: a sequence of abelian groups
linked by homomorphisms whose consecutive composites vanish.  The groups encode
objects in each dimension, while the homomorphisms encode boundaries.  This
chapter develops the algebra of chain complexes before any particular
topological chain model is chosen.

\section{Chain complexes}

\begin{definition}[Chain complex]
\label{def:viii16-chain-complex}
A chain complex \(C_\bullet\) of abelian groups is a sequence
\[
\cdots
\xrightarrow{\partial_{n+2}}
C_{n+1}
\xrightarrow{\partial_{n+1}}
C_n
\xrightarrow{\partial_n}
C_{n-1}
\xrightarrow{\partial_{n-1}}
\cdots
\]
such that
\[
\partial_n\partial_{n+1}=0
\]
for every \(n\).
\end{definition}

The homomorphisms \(\partial_n\) are boundary operators.

\section{Cycles and boundaries}

\begin{definition}[Cycles and boundaries]
\label{def:viii16-cycles-boundaries}
Define
\[
Z_n(C)=\ker\partial_n,
\qquad
B_n(C)=\operatorname{im}\partial_{n+1}.
\]
Elements of \(Z_n(C)\) are \(n\)-cycles and elements of \(B_n(C)\) are
\(n\)-boundaries.
\end{definition}

The identity \(\partial^2=0\) gives
\[
B_n(C)\subset Z_n(C).
\]

\section{Homology of a chain complex}

\begin{definition}[Homology group]
\label{def:viii16-homology}
The \(n\)th homology group of \(C_\bullet\) is
\[
H_n(C)=Z_n(C)/B_n(C).
\]
\end{definition}

At this stage homology is purely algebraic.  VIII/17 constructs chain
complexes from spaces and simplicial complexes.

\section{Concentrated complexes}

A group \(A\) may be viewed as a chain complex concentrated in degree \(k\):
\[
C_k=A,\qquad C_n=0\ \text{for }n\ne k.
\]

\begin{example}
For this complex,
\[
H_k(C)\cong A
\]
and every other homology group vanishes.
\end{example}

\section{Two-term complexes}

Consider
\[
0\longrightarrow A\xrightarrow{f}B\longrightarrow0,
\]
with \(A\) in degree \(1\) and \(B\) in degree \(0\).

\begin{proposition}[Two-term homology]
\label{prop:viii16-two-term}
\[
H_1\cong\ker f,
\qquad
H_0\cong B/\operatorname{im}f.
\]
\end{proposition}

Thus kernels and cokernels are already homology groups.

\section{Exact chain complexes}

\begin{definition}[Exactness]
\label{def:viii16-exact}
A chain complex is exact at \(C_n\) if
\[
\operatorname{im}\partial_{n+1}
=
\ker\partial_n.
\]
It is exact if this holds in every degree.
\end{definition}

\begin{proposition}[Exact iff acyclic]
\label{prop:viii16-exact-acyclic}
A chain complex is exact if and only if
\[
H_n(C)=0
\]
for all \(n\).
\end{proposition}

\section{Chain maps}

\begin{definition}[Chain map]
\label{def:viii16-chain-map}
A chain map
\[
f_\bullet:C_\bullet\to D_\bullet
\]
is a family of homomorphisms
\[
f_n:C_n\to D_n
\]
such that
\[
\partial^D_n f_n
=
f_{n-1}\partial^C_n
\]
for every \(n\).
\end{definition}

Equivalently, every boundary square commutes.

\section{Induced maps on homology}

\begin{theorem}[Homology is functorial on chain maps]
\label{thm:viii16-induced}
A chain map sends cycles to cycles and boundaries to boundaries, and therefore
induces
\[
H_n(f):H_n(C)\to H_n(D)
\]
by
\[
[z]\longmapsto[f_n(z)].
\]
\end{theorem}

\section{Identity and composition}

If \(f:C\to D\) and \(g:D\to E\) are chain maps, then \(g\circ f\) is a chain
map and
\[
H_n(g\circ f)=H_n(g)\circ H_n(f).
\]
Also
\[
H_n(\operatorname{id}_C)=\operatorname{id}_{H_n(C)}.
\]

\section{Chain isomorphisms}

\begin{definition}[Chain isomorphism]
\label{def:viii16-chain-iso}
A chain map \(f:C\to D\) is a chain isomorphism if every
\[
f_n:C_n\to D_n
\]
is an isomorphism.
\end{definition}

A chain isomorphism induces isomorphisms on every homology group.

\section{Subcomplexes}

\begin{definition}[Chain subcomplex]
\label{def:viii16-subcomplex}
A chain subcomplex \(A_\bullet\subset C_\bullet\) consists of subgroups
\[
A_n\subset C_n
\]
with
\[
\partial_C(A_n)\subset A_{n-1}.
\]
\end{definition}

The restricted maps form a chain complex.

\section{Quotient complexes}

If \(A_\bullet\subset C_\bullet\) is a chain subcomplex, define
\[
(C/A)_n=C_n/A_n.
\]

\begin{proposition}[Quotient boundary]
\label{prop:viii16-quotient}
The formula
\[
\partial[c]=[\partial c]
\]
is well defined and makes \(C_\bullet/A_\bullet\) a chain complex.
\end{proposition}

\section{Short exact sequences of complexes}

\begin{definition}[Short exact sequence of chain complexes]
\label{def:viii16-short-exact}
A sequence
\[
0\to A_\bullet\xrightarrow{i}B_\bullet\xrightarrow{p}C_\bullet\to0
\]
is short exact if
\[
0\to A_n\xrightarrow{i_n}B_n\xrightarrow{p_n}C_n\to0
\]
is exact for every \(n\), and all maps are chain maps.
\end{definition}

Such a sequence produces a long exact sequence in homology.  The construction
of the connecting homomorphism is developed in VIII/19.

\section{Direct sums}

Given chain complexes \(C^\alpha_\bullet\), define
\[
\left(\bigoplus_\alpha C^\alpha\right)_n
=
\bigoplus_\alpha C^\alpha_n
\]
with boundary acting componentwise.

\begin{proposition}[Homology and direct sums]
\label{prop:viii16-direct-sum}
For ordinary direct sums of abelian-group chain complexes,
\[
H_n\left(\bigoplus_\alpha C^\alpha\right)
\cong
\bigoplus_\alpha H_n(C^\alpha).
\]
\end{proposition}

\section{Augmented complexes}

An augmented chain complex continues one step beyond degree zero:
\[
C_1\xrightarrow{\partial_1}C_0
\xrightarrow{\epsilon}A\to0,
\]
with
\[
\epsilon\partial_1=0.
\]

For simplicial complexes, the augmentation
\[
\epsilon:C_0\to\mathbb Z
\]
sends every vertex to \(1\).  This leads to reduced homology in VIII/17.

\section{Reduced chain complexes}

The reduced degree-zero information is obtained by replacing the ordinary
degree-zero homology with the kernel of the augmentation.  For a connected
space this removes the universal \(\mathbb Z\) summand of \(H_0\).

\section{Graded viewpoint}

A chain complex may be written compactly as a graded group
\[
C_*=\bigoplus_n C_n
\]
equipped with a degree-\(-1\) map
\[
\partial:C_*\to C_*
\]
satisfying
\[
\partial^2=0.
\]

This notation becomes especially convenient for chain homotopies in VIII/22.

\section{Matrices for finite free complexes}

If every \(C_n\) is finitely generated free abelian, choose bases and represent
\[
\partial_n
\]
by an integer matrix \(D_n\).  The chain condition becomes
\[
D_{n-1}D_n=0.
\]

Homology can then be computed using kernels, images, and Smith normal form.

\section{Smith normal form preview}

For an integer matrix \(D\), unimodular row and column operations can put \(D\)
into diagonal form
\[
\operatorname{diag}(d_1,\dots,d_r,0,\dots,0),
\qquad
d_i\mid d_{i+1}.
\]

This is the standard computational tool for finitely generated homology groups.

\section{Euler characteristic preview}

For a finite chain complex of finitely generated free abelian groups, let
\[
c_n=\operatorname{rank}C_n.
\]
The alternating chain rank
\[
\sum_n(-1)^n c_n
\]
will later agree with the alternating sum of homology ranks.

\section{Bridge to topological chains}

VIII/15 provided oriented simplices and the alternating face formula.  In
VIII/17 those simplices become generators of chain groups and the face formula
becomes a boundary operator satisfying
\[
\partial^2=0.
\]
The present chapter supplies the algebra into which those constructions fit.


\section{Exercises}

\begin{exercise}\label{exr:viii16-01}
State the chain condition.
\end{exercise}
\begin{hint}
Compose consecutive boundaries.

% VIII-PEDAGOGY-HINT-v1.1 exr:viii16-01 append
Write the domains \(C_{n+1}\to C_n\to C_{n-1}\); the composite must be the zero homomorphism, not merely have a nontrivial kernel.
\end{hint}
\begin{solution}
\(\partial_n\partial_{n+1}=0\).
\end{solution}

\begin{exercise}\label{exr:viii16-02}
Define \(Z_n(C)\).
\end{exercise}
\begin{hint}
Take a kernel.

% VIII-PEDAGOGY-HINT-v1.1 exr:viii16-02 append
A cycle belongs to degree \(n\) and has zero outgoing boundary; identify the relevant kernel inside \(C_n\).
\end{hint}
\begin{solution}
\(Z_n(C)=\ker\partial_n\).
\end{solution}

\begin{exercise}\label{exr:viii16-03}
Define \(B_n(C)\).
\end{exercise}
\begin{hint}
Take the preceding image.

% VIII-PEDAGOGY-HINT-v1.1 exr:viii16-03 append
A degree-\(n\) boundary comes from degree \(n+1\), so use the incoming differential rather than the outgoing one.
\end{hint}
\begin{solution}
\(B_n(C)=\operatorname{im}\partial_{n+1}\).
\end{solution}

\begin{exercise}\label{exr:viii16-04}
Why is \(B_n\subset Z_n\)?
\end{exercise}
\begin{hint}
Use \(\partial^2=0\).

% VIII-PEDAGOGY-HINT-v1.1 exr:viii16-04 append
Write an arbitrary boundary as \(b=\partial_{n+1}c\), then apply \(\partial_n\) and use the chain condition.
\end{hint}
\begin{solution}
If \(b=\partial_{n+1}c\), then \(\partial_nb=\partial_n\partial_{n+1}c=0\).
\end{solution}

\begin{exercise}\label{exr:viii16-05}
Define \(H_n(C)\).
\end{exercise}
\begin{hint}
Cycles modulo boundaries.

% VIII-PEDAGOGY-HINT-v1.1 exr:viii16-05 append
First justify that boundaries form a subgroup of cycles; two cycle representatives define the same class exactly when their difference is a boundary.
\end{hint}
\begin{solution}
\(H_n(C)=Z_n(C)/B_n(C)\).
\end{solution}

\begin{exercise}\label{exr:viii16-06}
What is the homology of a complex concentrated in degree \(k\) at \(A\)?
\end{exercise}
\begin{hint}
All boundaries vanish.

% VIII-PEDAGOGY-HINT-v1.1 exr:viii16-06 append
At the sole nonzero degree, the outgoing differential has full kernel and the incoming differential has zero image; inspect the other degrees separately.
\end{hint}
\begin{solution}
\(H_k\cong A\), all other \(H_n=0\).
\end{solution}

\begin{exercise}\label{exr:viii16-07}
For \(0\to A\xrightarrow f B\to0\), what is \(H_1\)?
\end{exercise}
\begin{hint}
Take the kernel of \(f\).

% VIII-PEDAGOGY-HINT-v1.1 exr:viii16-07 append
Place \(A\) in degree one and \(B\) in degree zero; there is no incoming degree-two boundary to quotient out.
\end{hint}
\begin{solution}
\(\ker f\).
\end{solution}

\begin{exercise}\label{exr:viii16-08}
For the same complex, what is \(H_0\)?
\end{exercise}
\begin{hint}
Take the cokernel.

% VIII-PEDAGOGY-HINT-v1.1 exr:viii16-08 append
The outgoing map from \(B\) is zero, so all of \(B\) consists of cycles; divide by the image arriving from \(A\).
\end{hint}
\begin{solution}
\(B/\operatorname{im}f\).
\end{solution}

\begin{exercise}\label{exr:viii16-09}
Define exactness at \(C_n\).
\end{exercise}
\begin{hint}
Image equals kernel.

% VIII-PEDAGOGY-HINT-v1.1 exr:viii16-09 append
The chain condition gives only image contained in kernel; exactness requires the reverse containment as well.
\end{hint}
\begin{solution}
\(\operatorname{im}\partial_{n+1}=\ker\partial_n\).
\end{solution}

\begin{exercise}\label{exr:viii16-10}
When is a chain complex acyclic?
\end{exercise}
\begin{hint}
All homology vanishes.

% VIII-PEDAGOGY-HINT-v1.1 exr:viii16-10 append
Apply the cycles-modulo-boundaries definition in every degree, distinguishing ordinary acyclicity from reduced acyclicity of a topological space.
\end{hint}
\begin{solution}
When \(H_n(C)=0\) for all \(n\).
\end{solution}

\begin{exercise}\label{exr:viii16-11}
Define a chain map.
\end{exercise}
\begin{hint}
Boundary squares commute.

% VIII-PEDAGOGY-HINT-v1.1 exr:viii16-11 append
For each \(n\), compare the two maps \(C_n\to D_{n-1}\) obtained by applying the differential before or after \(f\).
\end{hint}
\begin{solution}
\(\partial^Df_n=f_{n-1}\partial^C\).
\end{solution}

\begin{exercise}\label{exr:viii16-12}
Why does a chain map send cycles to cycles?
\end{exercise}
\begin{hint}
Apply the commuting square.

% VIII-PEDAGOGY-HINT-v1.1 exr:viii16-12 append
Substitute a cycle \(z\) into \(\partial^D f_n=f_{n-1}\partial^C\); its zero boundary should remain zero after applying \(f\).
\end{hint}
\begin{solution}
If \(\partial z=0\), then \(\partial f(z)=f(\partial z)=0\).
\end{solution}

\begin{exercise}\label{exr:viii16-13}
Why does a chain map send boundaries to boundaries?
\end{exercise}
\begin{hint}
Move the boundary through \(f\).

% VIII-PEDAGOGY-HINT-v1.1 exr:viii16-13 append
Write the input as \(\partial^C_{n+1}c\) and use the chain-map equation in degree \(n+1\) to exhibit an explicit boundary preimage.
\end{hint}
\begin{solution}
\(f(\partial c)=\partial f(c)\).
\end{solution}

\begin{exercise}\label{exr:viii16-14}
What map does a chain map induce?
\end{exercise}
\begin{hint}
Pass to quotient classes.

% VIII-PEDAGOGY-HINT-v1.1 exr:viii16-14 append
Define the image of a cycle class by \([z]\mapsto[f_n(z)]\), then check that changing \(z\) by a boundary does not change the output class.
\end{hint}
\begin{solution}
\(H_n(f):H_n(C)\to H_n(D)\).
\end{solution}

\begin{exercise}\label{exr:viii16-15}
What is a chain isomorphism?
\end{exercise}
\begin{hint}
Degreewise isomorphism.

% VIII-PEDAGOGY-HINT-v1.1 exr:viii16-15 append
Invert the degreewise isomorphisms in the commuting differential square to verify that their inverses also form a chain map.
\end{hint}
\begin{solution}
A chain map whose every component \(f_n\) is an isomorphism.
\end{solution}

\begin{exercise}\label{exr:viii16-16}
Define a chain subcomplex.
\end{exercise}
\begin{hint}
Boundary preserves the subgroups.

% VIII-PEDAGOGY-HINT-v1.1 exr:viii16-16 append
Containment \(A_n\subset C_n\) alone is insufficient: apply the ambient differential and require its image to lie in \(A_{n-1}\).
\end{hint}
\begin{solution}
Subgroups \(A_n\subset C_n\) with \(\partial(A_n)\subset A_{n-1}\).
\end{solution}

\begin{exercise}\label{exr:viii16-17}
How is the quotient boundary defined?
\end{exercise}
\begin{hint}
Apply boundary before taking the coset.

% VIII-PEDAGOGY-HINT-v1.1 exr:viii16-17 append
Replace a representative \(c\) by \(c+a\) with \(a\in A_n\); boundary stability of the subcomplex makes both images give the same quotient class.
\end{hint}
\begin{solution}
\(\partial[c]=[\partial c]\).
\end{solution}

\begin{exercise}\label{exr:viii16-18}
Define short exactness for chain complexes.
\end{exercise}
\begin{hint}
Require exactness degreewise.

% VIII-PEDAGOGY-HINT-v1.1 exr:viii16-18 append
Require chain maps as well as exactness of the underlying groups in each degree; these commuting differentials enable the homology construction.
\end{hint}
\begin{solution}
Every \(0\to A_n\to B_n\to C_n\to0\) is exact.
\end{solution}

\begin{exercise}\label{exr:viii16-19}
What does a short exact sequence of complexes later produce?
\end{exercise}
\begin{hint}
Recall VIII/19.

% VIII-PEDAGOGY-HINT-v1.1 exr:viii16-19 append
Start with a cycle in the quotient complex, lift it to the middle complex, and take its boundary to locate the degree-lowering connecting map.
\end{hint}
\begin{solution}
A long exact sequence in homology.
\end{solution}

\begin{exercise}\label{exr:viii16-20}
How is the boundary on a direct sum defined?
\end{exercise}
\begin{hint}
Componentwise.

% VIII-PEDAGOGY-HINT-v1.1 exr:viii16-20 append
Elements of a direct sum have finite support; applying each differential separately preserves that finite-support condition.
\end{hint}
\begin{solution}
\(\partial((c_\alpha))=(\partial c_\alpha)\).
\end{solution}

\begin{exercise}\label{exr:viii16-21}
What is an augmentation?
\end{exercise}
\begin{hint}
Continue beyond \(C_0\).

% VIII-PEDAGOGY-HINT-v1.1 exr:viii16-21 append
Treat the augmentation as an extra differential leaving degree zero and verify its composite with \(\partial_1\) vanishes.
\end{hint}
\begin{solution}
A map \(\epsilon:C_0\to A\) satisfying \(\epsilon\partial_1=0\).
\end{solution}

\begin{exercise}\label{exr:viii16-22}
What does the standard simplicial augmentation send a vertex to?
\end{exercise}
\begin{hint}
Every vertex contributes one component.

% VIII-PEDAGOGY-HINT-v1.1 exr:viii16-22 append
Evaluate augmentation on the edge boundary \([v_1]-[v_0]\); sending both vertices to one makes the result zero.
\end{hint}
\begin{solution}
\(1\in\mathbb Z\).
\end{solution}

\begin{exercise}\label{exr:viii16-23}
What matrix identity expresses the chain condition?
\end{exercise}
\begin{hint}
Choose bases.

% VIII-PEDAGOGY-HINT-v1.1 exr:viii16-23 append
Check matrix sizes against \(C_n\to C_{n-1}\to C_{n-2}\), then multiply in that order using a consistent column-vector convention.
\end{hint}
\begin{solution}
\(D_{n-1}D_n=0\).
\end{solution}

\begin{exercise}\label{exr:viii16-24}
What is the bridge to VIII/17?
\end{exercise}
\begin{hint}
Use oriented simplices.

% VIII-PEDAGOGY-HINT-v1.1 exr:viii16-24 append
Apply the alternating face operator twice to one generator; pairing the two deletion orders for each pair of vertices supplies the chain condition.
\end{hint}
\begin{solution}
VIII/17 turns simplices into chain generators and the alternating face sum into the boundary operator.
\end{solution}


\section{Solved problem dossiers}

\begin{problem}[Zero differential]\label{prob:viii16-01}
Let \(C_n=A_n\) with every boundary zero. Compute \(H_n(C)\).
\end{problem}
\begin{solution}
Every element is a cycle and no nonzero element is a boundary, so \(H_n(C)\cong A_n\) for every \(n\).
\end{solution}

\begin{problem}[Identity differential warning]\label{prob:viii16-02}
Can a nonzero chain complex have every differential equal to the identity?
\end{problem}
\begin{solution}
Not in consecutive nonzero degrees, because the chain condition would give \(\operatorname{id}\circ\operatorname{id}=0\), impossible on a nonzero group.
\end{solution}

\begin{problem}[Multiplication-by-\(m\) complex]\label{prob:viii16-03}
Compute the homology of \(0\to\mathbb Z\xrightarrow{\times m}\mathbb Z\to0\) for \(m\ne0\).
\end{problem}
\begin{solution}
The map is injective, so \(H_1=0\). Its image is \(m\mathbb Z\), so \(H_0\cong\mathbb Z/m\).
\end{solution}

\begin{problem}[Zero multiplication]\label{prob:viii16-04}
Repeat the previous computation for \(m=0\).
\end{problem}
\begin{solution}
The boundary is zero, hence \(H_1\cong\mathbb Z\) and \(H_0\cong\mathbb Z\).
\end{solution}

\begin{problem}[Exactness test]\label{prob:viii16-05}
Show \(0\to\mathbb Z\xrightarrow{\operatorname{id}}\mathbb Z\to0\) is exact.
\end{problem}
\begin{solution}
The identity has zero kernel and full image. Thus the degree-one and degree-zero homology groups both vanish.
\end{solution}

\begin{problem}[Cycle-to-cycle proof]\label{prob:viii16-06}
Prove a chain map preserves cycles.
\end{problem}
\begin{solution}
For \(z\in Z_n(C)\), \(\partial^Df_n(z)=f_{n-1}\partial^C(z)=0\), so \(f_n(z)\in Z_n(D)\).
\end{solution}

\begin{problem}[Boundary-to-boundary proof]\label{prob:viii16-07}
Prove a chain map preserves boundaries.
\end{problem}
\begin{solution}
If \(b=\partial^C_{n+1}c\), then \(f_n(b)=f_n\partial^C c=\partial^D f_{n+1}(c)\), so it is a boundary.
\end{solution}

\begin{problem}[Well-defined induced map]\label{prob:viii16-08}
Why does \([z]\mapsto[f(z)]\) not depend on the chosen cycle representative?
\end{problem}
\begin{solution}
If \(z-z'=\partial c\), then \(f(z)-f(z')=f(\partial c)=\partial f(c)\), so the images differ by a boundary.
\end{solution}

\begin{problem}[Composition on homology]\label{prob:viii16-09}
Verify \(H_n(gf)=H_n(g)H_n(f)\).
\end{problem}
\begin{solution}
On a class \([z]\), both sides equal \([g_n(f_n(z))]\).
\end{solution}

\begin{problem}[Quotient differential]\label{prob:viii16-10}
Prove \(\partial[c]=[\partial c]\) is well defined on \(C/A\).
\end{problem}
\begin{solution}
If \(c-c'\in A_n\), then \(\partial(c-c')\in A_{n-1}\), so \([\partial c]=[\partial c']\).
\end{solution}

\begin{problem}[Quotient chain condition]\label{prob:viii16-11}
Verify the quotient differential squares to zero.
\end{problem}
\begin{solution}
\(\partial^2[c]=[\partial^2c]=[0]\).
\end{solution}

\begin{problem}[Direct sum homology]\label{prob:viii16-12}
Why do cycles and boundaries of a direct sum split componentwise?
\end{problem}
\begin{solution}
The direct-sum differential acts componentwise, so kernels and images are direct sums of the component kernels and images, giving the homology direct-sum formula.
\end{solution}

\begin{problem}[Augmentation of a connected graph]\label{prob:viii16-13}
Why does sending each vertex generator to \(1\) annihilate every edge boundary?
\end{problem}
\begin{solution}
An oriented edge has boundary \(v_1-v_0\), and the augmentation sends this to \(1-1=0\).
\end{solution}

\begin{problem}[Matrix chain condition]\label{prob:viii16-14}
Let \(D_2=\begin{pmatrix}1\\1\end{pmatrix}\) and \(D_1=\begin{pmatrix}1&-1\end{pmatrix}\). Verify the chain condition.
\end{problem}
\begin{solution}
\(D_1D_2=1-1=0\).
\end{solution}

\begin{problem}[Matrix homology]\label{prob:viii16-15}
For the preceding complex \(0\to\mathbb Z\xrightarrow{D_2}\mathbb Z^2\xrightarrow{D_1}\mathbb Z\to0\), compute the homology.
\end{problem}
\begin{solution}
\(D_2\) is injective, so \(H_2=0\). The kernel of \(D_1\) is generated by \((1,1)\), exactly the image of \(D_2\), so \(H_1=0\). \(D_1\) is surjective, so \(H_0=0\).
\end{solution}

\begin{problem}[Nonexact matrix example]\label{prob:viii16-16}
Replace \(D_2\) above by multiplication of the generator by \((2,2)\). Compute \(H_1\).
\end{problem}
\begin{solution}
\(\ker D_1=\mathbb Z(1,1)\) and \(\operatorname{im}D_2=2\mathbb Z(1,1)\), so \(H_1\cong\mathbb Z/2\).
\end{solution}

\begin{problem}[Chain isomorphism]\label{prob:viii16-17}
Why must a chain isomorphism induce homology isomorphisms?
\end{problem}
\begin{solution}
Its inverse is also a chain map, and functoriality gives mutually inverse induced maps on homology.
\end{solution}

\begin{problem}[Exact iff zero homology]\label{prob:viii16-18}
Prove the equivalence.
\end{problem}
\begin{solution}
Exactness says \(Z_n=B_n\) in every degree, which is equivalent to the quotient \(Z_n/B_n\) being zero.
\end{solution}

\begin{problem}[Short exact quotient example]\label{prob:viii16-19}
Given a subcomplex \(A\subset C\), write the canonical short exact sequence.
\end{problem}
\begin{solution}
\(0\to A_\bullet\to C_\bullet\to(C/A)_\bullet\to0\), degreewise the standard subgroup-quotient sequence.
\end{solution}

\begin{problem}[Smith-form preview]\label{prob:viii16-20}
What does the complex \(0\to\mathbb Z\xrightarrow{\times12}\mathbb Z\to0\) predict for torsion?
\end{problem}
\begin{solution}
The single Smith invariant is \(12\), and the degree-zero homology is \(\mathbb Z/12\).
\end{solution}

\begin{problem}[Euler rank check]\label{prob:viii16-21}
For a finite exact complex of finite-rank free groups, what alternating rank should one expect?
\end{problem}
\begin{solution}
Zero. Exactness lets ranks telescope across images and kernels; later the Euler--Poincare formula formalizes this.
\end{solution}

\begin{problem}[Reduced degree zero]\label{prob:viii16-22}
Why is reduced homology useful for connected spaces?
\end{problem}
\begin{solution}
Ordinary \(H_0\) always contains the universal \(\mathbb Z\) component recording connectedness. The augmentation removes that component, making reduced \(H_0\) vanish for connected spaces.
\end{solution}
